% ECONOMICS_PROBLEM: {"id": "C73-40", "title": "Nonlinear stability of periodic MFG solutions", "jel_code": "C73", "source_id": "OP-0211"}
% !TeX program = xelatex
\documentclass[11pt,a4paper]{article}
\usepackage[margin=23mm]{geometry}
\usepackage{fontspec}
\setmainfont{Latin Modern Roman}
\usepackage{amsmath,amssymb,mathtools}
\usepackage{xcolor,enumitem,booktabs,longtable,array,xurl,hyperref}
\definecolor{muted}{HTML}{53636C}
\hypersetup{colorlinks=true,urlcolor=blue,linkcolor=blue}
\setlength{\parindent}{0pt}
\setlength{\parskip}{5pt}
\newcommand{\R}{\mathbb R}
\newcommand{\E}{\mathbb E}
\newcommand{\N}{\mathbb N}
\newcommand{\ind}{\mathbf 1}
\newcommand{\Var}{\operatorname{Var}}
\newcommand{\Cov}{\operatorname{Cov}}
\newcommand{\supp}{\operatorname{supp}}
\DeclareMathOperator*{\argmin}{arg\,min}
\DeclareMathOperator*{\argmax}{arg\,max}
\newcommand{\entrymeta}[3]{{\sffamily\small\color{muted}\textbf{\texttt{#1}}\quad|\quad #2\par #3\par}}
\newcommand{\Problemref}[1]{\texttt{#1}}
\newcommand{\JELref}[2]{\texttt{#2} (JEL \texttt{#1})}
\begin{document}
\section*{Nonlinear stability of periodic MFG solutions}
\textbf{Problem ID:} \texttt{C73-40}\quad \textbf{JEL:} \texttt{C73}\par
% BEGIN ECONOMICS STATEMENT
\label{problem:OP-0211}
% GAME_THEORY_PROBLEM: {"local_id": "GT-081", "original_number": 81, "kind": "R", "entry_kind": "statement_only", "published": false, "review_required": true, "category": "game_theory"}
\label{game:GT-081}
{\sffamily\small\color{muted}
{\small\textbf{GT-081} | Mean-field games | R: an explicit boundary-value stability formulation\par}
\par}
\subsection*{Definitions and mathematical statement}
Fix $\nu>0$ and a smooth local coupling $f:\mathbb T^d\times(0,\infty)\to\mathbb R$. Suppose a classical triple $(v_*,m_*,\lambda_*)$ solves on $\mathbb R\times\mathbb T^d$
\[
 \begin{cases}
 \lambda_*-\partial_t v_*-\nu\Delta v_*+\tfrac12|Dv_*|^2=f(x,m_*),\\
 \partial_t m_*-\nu\Delta m_*-\operatorname{div}(m_*Dv_*)=0,\\
 v_*(t+P)=v_*(t),\quad m_*(t+P)=m_*(t),\quad
 m_*>0,\quad\int m_*(t)=1,
 \end{cases}
\]
for some $P>0$, with $(Dv_*,m_*)$ nonconstant in time. The associated unnormalized value is $u_*(t,x)=v_*(t,x)-\lambda_*t$; $u_*$ itself need not be periodic. Define its orbit
\[
 \mathcal O_* =\{(Dv_*(\theta),m_*(\theta)):0\leq\theta<P\},
 \quad
 d_*((p,m),\mathcal O_*)=
 \inf_{0\leq\theta<P}(\|p-Dv_*(\theta)\|_\infty+\|m-m_*(\theta)\|_1).
\]
Fix constants $M,M_0<\infty$ and $\zeta>0$ containing the reference solution in their bounds. The admissible solution class has $m\geq\zeta$ and $\|m\|_{C^{1,2}}+\|Du\|_\infty+\|D^2u\|_\infty\leq M$; smooth boundary data have $m_0\geq\zeta$, $\int m_0=1$, and $\|m_0\|_{C^2}+\|Dh\|_{C^1}\leq M_0$. The $C^{1,2}$ norm controls the function, its first time derivative, and its spatial derivatives through order two. One precise editorial stability notion is \emph{uniform boundary-value orbital shadowing}. For every $\varepsilon>0$, require a $\delta>0$, independent of $T\geq1$ and phase $\theta$, such that whenever these boundary data obey
\[
 \|m_0-m_*(\theta)\|_1+\|Dh-Dv_*(T+\theta)\|_\infty<\delta,
\]
every admissible classical solution of
\[
 -u_t-\nu\Delta u+\tfrac12|Du|^2=f(x,m),\qquad
 m_t-\nu\Delta m-\operatorname{div}(mDu)=0,
 \quad m(0)=m_0,\quad u(T)=h
\]
satisfies $\sup_{0\leq t\leq T}d_*((Du(t),m(t)),\mathcal O_*)<\varepsilon$. Nonemptiness of the admissible solution class for all such perturbed boundary data is also required; vacuous bounds on an empty class do not establish stability.

Determine which periodic equilibria satisfy this property, or a stated weaker nonlinear stability property with an explicit equilibrium-selection rule. This is a boundary-value comparison problem; specifying arbitrary $m(0),u(0)$ does not define an MFG evolution semigroup.
\subsection*{Short English statement}
Do equilibria near a periodic solution remain near its orbit when their compatible initial and terminal data are perturbed?
\subsection*{Variants and scope boundaries}
Linearized spectral stability, orbital shadowing, and asymptotic attraction are distinct. Terminal costs depending on the terminal distribution require control of that additional coupling. A finite-horizon oscillation is not automatically an entire periodic solution.
\subsection*{Source relationship and status}
The chapter identifies stability as unknown for its examples. The chosen stability definition is editorial; no all-periodic-solutions stability theorem is asserted.
\subsection*{References}
\begingroup\footnotesize\raggedright
{\footnotesize\raggedright
\textbf{[S1]} Cardaliaguet and Delarue. \emph{Selected topics in mean field games}, ICM2022, vol.~5; DOI: 10.4171/ICM2022/17. \textit{Locator:} Section~3.3, ``Periodic solutions,'' p.~3680.\par
\url{https://ems.press/content/book-chapter-files/33265}\par\smallskip
\textbf{[S2]} Marco Cirant and Levon Nurbekyan (2018). \emph{The variational structure and time-periodic solutions for mean-field games systems}. Minimax Theory and its Applications 3, 227--260.\par
\textit{Locators:} arXiv version, Section~6, Theorem~6.1, pp.~24--25, actual periodic examples; Remark~6.4, p.~28, limited qualitative information.\par
\url{https://arxiv.org/pdf/1804.08943v1}\par}
\par\endgroup

\subsection*{Common conventions: Source mathematical conventions}

\textbf{Finite games.} For a finite set $A$, $\Delta(A)$ is its probability
simplex. Players choose independent mixed strategies in Nash equilibrium;
correlation is allowed only when explicitly part of the solution concept.
In a game with payoff functions $u_i$ and strategy profile $x$, an additive
$\varepsilon$ Nash equilibrium satisfies
\[
 u_i(x)\geq u_i(x_i',x_{-i})-\varepsilon
 \quad\text{for every player $i$ and every admissible deviation $x_i'$}.
\]
For numerical approximation, payoffs are normalized as locally specified.
An $\varepsilon$ payoff residual is distinct from distance at most
$\varepsilon$ to the set of exact equilibria. Point convergence, convergence
of empirical play, and convergence in distance to an equilibrium set are
also distinct.

\textbf{Computational representations.} Unless stated otherwise, a complexity
question takes rational data with binary encoding; polynomial time is measured
in total bit length. A fixed-accuracy polynomial algorithm is not necessarily
polynomial in $1/\varepsilon$, and neither is necessarily polynomial in
$\log(1/\varepsilon)$. Succinct games and value, demand, or sampling oracles
must declare their interfaces. Exact real solutions require an explicit
representation; an unspecified list of arbitrary real numbers is not a
finite computational output. A claimed algorithmic barrier is meaningful
only for its specified model and reductions.

\textbf{Dynamic games.} Histories, information, observation, and allowable
randomization are part of the model. A stationary strategy depends only on
the current state; a positional strategy in a deterministic graph game
chooses one action at each controlled vertex. Nash equilibrium at an initial
state and equilibrium after every history are different requirements.
An expectation of a pathwise $\liminf$ payoff, a $\liminf$ of expected averages,
a limiting expected average, and a uniform finite-horizon equilibrium are
different criteria. None is silently substituted for another.

\textbf{Allocation and mechanisms.} A complete allocation assigns every item
exactly once unless otherwise stated. Zero-valued goods, negative values,
capacity constraints, feasibility, lotteries, and copies of goods affect
fairness definitions and are specified locally. Dominant-strategy incentive
compatibility, Bayesian incentive compatibility, universal truthfulness,
and truthfulness in expectation impose different inequalities. Whenever
an objective ratio has zero denominator, its convention is stated or those
instances are excluded.

\textbf{Infinite-dimensional models.} Expectations, integrals, stochastic
equations, and optimization objectives require their local regularity,
integrability, filtrations, and admissible domains. A backward--forward
mean-field system is a boundary-value problem; it is not an initial-value
semiflow without an additional construction. Long-time, large-population,
and vanishing-noise limits require an order of limits and a convergence
topology. If a minimum need not be attained, the objective is its infimum
or supremum and the approximation requirement is stated separately.
% END ECONOMICS STATEMENT
\end{document}
